\section{Overview of the full development}\label{ch:intro}
\sectionbrief
{one function $u(t,x)$, known exactly (no noise) on an open region $\Omega\subset\R^2$, which solves a linear equation $P(\partial_t,\partial_x)u=0$ with constant coefficients. The operator $P$ is \emph{not} given.}
{to state precisely what we want to learn from $u$ and to give a map of the answers before any detail.}
{frames the three questions Q1--Q3. Sections~\ref{f:ch:ident}--\ref{f:ch:data} answer Q1, Sections~\ref{f:ch:decomp}--\ref{f:ch:regular} answer Q2, Section~\ref{f:ch:interaction} answers Q3. Section~\ref{f:ch:nonlin} extends all three to nonlinear equations and Section~\ref{ch:algo} assembles them into an algorithm. Part~II (Sections~\ref{ch:systems}--\ref{f:ch:status}) goes beyond constant coefficients and states the limits of the whole project.}

\subsection{The question}

Suppose we observe a physical field $u(t,x)$: a temperature, a wave, a concentration. We know $u$ exactly on some region $\Omega$ of the $(t,x)$ plane, and we know that it obeys some linear law. We do not know which law. Three questions arise naturally.

\begin{enumerate}[label=\textbf{Q\arabic*.},leftmargin=1.2cm]
\item \textbf{Identification.} Can the operator $P$ be reconstructed from $u$? When is the answer unique?
\item \textbf{Decomposition.} If $P$ is built from pieces, can $u$ be split into the parts that each piece generates?
\item \textbf{Interaction.} If the pieces are the \emph{terms} of the operator (diffusion, advection, dispersion, \dots), which part of $u$ is due to no single term but to their interplay? Can it be computed exactly and cheaply?
\end{enumerate}

The setting is deliberately analytic. There is no statistical model, no learned weights and no noise. The aim is an \emph{exact} procedure, together with a precise statement of what is impossible.

\begin{example}[{Three solutions of the heat equation}]
Take the heat equation $u_t=Du_{xx}$ with $D>0$. Consider three of its solutions.
\begin{itemize}
\item A single Fourier mode:
\[
u_1=e^{-D\kappa^2t}\cos(\kappa x).
\]
\item A sum of two modes:
\[
u_2=e^{-Dt}\cos x+e^{-4Dt}\cos 2x.
\]
\item The heat kernel, for $t>0$:
\[
u_3=\frac{1}{\sqrt{t}}\;e^{-x^2/(4Dt)}.
\]
\end{itemize}
All three solve the same equation. We will see that only $u_3$ determines it. The function $u_1$ also solves $u_{xx}=-\kappa^2u$ and $u_t=-D\kappa^2u$, and from $u_1$ alone we cannot tell which of these laws is ``the'' law. The difference is that $u_1$ and $u_2$ are finite sums of exponentials, while $u_3$ is not (Section~\ref{f:ch:ident}, Example~\ref{ex:heat3}).
\end{example}

\subsection{Overview of the answers}

Figure~\ref{fig:pipeline} shows the logical chain, read from left to right. Each box is a task. Each arrow is a theorem that allows the next task to be carried out, under its own hypotheses.

\begin{figure}[!htb]
\centering
\begin{tikzpicture}[node distance=7mm,
  box/.style={draw=cA,fill=cA!8,rounded corners=2pt,minimum height=16mm,text width=2.75cm,align=center,font=\small\sffamily},
  arr/.style={-{Stealth[length=2.4mm]},thick,cA}]
\node[box] (u) {exact solution\\$u$ on $\Omega$\\[2pt]{\scriptsize\textcolor{cG}{data}}};
\node[box,right=of u] (id) {\textbf{Identify}\\ $P$ and $G_u$\\[2pt]{\scriptsize\textcolor{cG}{Ch.~\ref{f:ch:ident}--\ref{f:ch:data}}}};
\node[box,right=of id] (fa) {\textbf{Factor}\\ $u=\sum_ju_j$\\[2pt]{\scriptsize\textcolor{cG}{Ch.~\ref{f:ch:decomp}--\ref{f:ch:regular}}}};
\node[box,right=of fa] (te) {\textbf{Terms}\\ log-symbol, $\Eint$\\[2pt]{\scriptsize\textcolor{cG}{Ch.~\ref{f:ch:interaction}}}};
\node[box,below=10mm of te,fill=cB!8,draw=cB] (nl) {\textbf{Nonlinear}\\ dictionary, flows\\[2pt]{\scriptsize\textcolor{cG}{Ch.~\ref{f:ch:nonlin}}}};
\node[box,below=10mm of id,fill=cC!8,draw=cC] (al) {\textbf{Algorithm}\\ Steps 0--7\\[2pt]{\scriptsize\textcolor{cG}{Ch.~\ref{ch:algo}}}};
\draw[arr] (u)--(id); \draw[arr] (id)--(fa); \draw[arr] (fa)--(te); \draw[arr] (te)--(nl);
\draw[arr,cC,dashed] (nl)--(al);
\end{tikzpicture}
\caption{Logical structure of the method. The three linear stages (blue) are followed by the nonlinear extension (red). The algorithm (green) chains all of them. A step can only be trusted under the hypotheses of the theorem that precedes it.}\label{fig:pipeline}
\end{figure}

Table~\ref{f:tab:results} lists the main statements with their status. Every row points to the chapter where the statement is made precise.

\begin{table}[!htb]
\centering\small
\begin{tabular}{@{}L{2.8cm}L{8.6cm}L{3.0cm}@{}}
\toprule
Topic & Result & Status\\\midrule
Identification & $\Ann(u)=(P)$ iff no $w_i=(P/P_i)(\partial)u$ is an exponential-polynomial (Ch.~\ref{f:ch:ident}) & \proved\\
Spectral form & $\Ann(u)=(P)$ iff the point spectrum is empty and $G_u\sim P$ (Ch.~\ref{f:ch:ident}) & \proved\\
Finite data & a jet of order $2d-1$ suffices for exponential-polynomials; no finite certificate otherwise; black box impossible; decidable with a bound (Ch.~\ref{f:ch:data}) & \proved\\
Decomposition & $\ker P=\sum_j\ker Q_j$ on simply connected $\Omega$, ambiguity $\delta(P)$ (Ch.~\ref{f:ch:decomp}) & \proved\\
Terms & splitting per additive term is false (Ch.~\ref{f:ch:decomp}) & \proved\\
Regularity & components in $H^m$ iff the sum of kernels is closed ($\gamma_N>0$) and dense (Ch.~\ref{f:ch:regular}) & \proved\\
Interaction & closed product formula; $O(N)$ per value, $\Theta(2^N)$ per table (Ch.~\ref{f:ch:interaction}) & \proved\\
Stability & $u_0\mapsto\Eint$ is not continuous in $H^m$; no stable regularisation (Ch.~\ref{f:ch:interaction}) & \proved\\
Nonlinear & Burgers: jet, closed form before the shock, $\Theta(\nu)$ and $\nu|\log\nu|$ after (Ch.~\ref{f:ch:nonlin}) & \proved{} / \modulo{Kru\v zkov}\\
Dictionary & identification modulo differential consequences; KdV degree $(N+2)2^{N-1}$ (Ch.~\ref{f:ch:nonlin}) & \proved{} / \modulo{Novikov}\\\bottomrule
\end{tabular}
\caption{Main results, in the order in which they appear.}\label{f:tab:results}
\end{table}

\subsection{Notation at a glance}

Table~\ref{tab:notation} collects the symbols used throughout, with the place where each one is defined. Some letters have different meanings in different chapters; the caption lists them.

\begin{table}[!htb]
\centering\small
\begin{tabular}{@{}L{3.1cm}L{8.9cm}L{2.5cm}@{}}
\toprule
Symbol & Meaning & Defined in\\\midrule
$(t,x)$, $\Omega$ & time and space; the open region where $u$ is known & Ch.~\ref{ch:intro}\\
$\partial=(\partial_t,\partial_x)$ & partial derivatives & Ch.~\ref{f:ch:algebra}\\
$\ring=\C[\tau,\xi]$ & ring of constant-coefficient operators, $\tau\leftrightarrow\partial_t$, $\xi\leftrightarrow\partial_x$ & Def.~\ref{f:def:ring}\\
$P$, $P_i$, $e_i$ & the operator (its symbol), its irreducible factors, their multiplicities & Ch.~\ref{f:ch:algebra}\\
$Q_j=P_j^{m_j}$ & pairwise coprime factors used for decomposition & Ch.~\ref{f:ch:decomp}\\
$\Var(P)$ & characteristic variety: zeros of the symbol in $\C^2$ & Def.~\ref{f:def:ring}\\
$\Ann(u)$ & the ideal of all operators that annihilate $u$ & Def.~\ref{f:def:ann}\\
$G_u$, $J_u$ & $G_u=\Gcd\Ann(u)$ and the cofactor ideal: $\Ann(u)=G_u\,J_u$ & Lemma~\ref{f:lem:gcd}\\
$\EP$ & exponential-polynomials & Def.~\ref{f:def:ep}\\
$\sigma_p(u)$ & point spectrum $\Var(J_u)$ & Ch.~\ref{f:ch:ident}\\
$\delta(P)$ & dimension of the ambiguity of the factor decomposition & Ch.~\ref{f:ch:decomp}\\
$c_F$, $\gamma_N$ & Friedrichs cosine; conditioning number of the sum map & Ch.~\ref{f:ch:regular}\\
$A_n$, $S_n(t)$, $\Delta_n$ & terms of an evolution, their flows, increments $S_n-\Id$ & Ch.~\ref{f:ch:interaction}\\
$\Eint$ & interaction term & Eq.~\eqref{f:eq:eint-def}\\
$T_n$, $V$, $I(u)$ & nonlinear terms; dictionary span; ideal of relations of $u$ & Ch.~\ref{f:ch:nonlin}\\
\bottomrule
\end{tabular}
\caption{Notation guide. Overloaded letters: $S$ is a flow in Section~\ref{f:ch:interaction} and the sum map in Section~\ref{f:ch:regular}; $D$ is $d/dx$ in Section~\ref{sec:ode} and a diffusion constant elsewhere; $J$ denotes ideals.}\label{tab:notation}
\end{table}

\subsection{Roadmap}

The chapters build on each other in the following order.
\begin{itemize}
\item \textbf{Section~\ref{ch:primer}} is a primer. It explains ideals, annihilators, exponential-polynomials, the Chinese remainder theorem and the companion matrix, starting from integers and one-variable ODEs. A reader who already knows these notions can skip it.
\item \textbf{Section~\ref{f:ch:algebra}} fixes the algebraic language: symbols, the annihilator ideal, its gcd splitting and the finite-dimension test for exponential-polynomials. It ends with the one-variable prototype.
\item \textbf{Sections~\ref{f:ch:ident} and~\ref{f:ch:data}} answer Q1. The first gives the exact criterion and the intrinsic operator $G_u$. The second asks what can be decided from finite data.
\item \textbf{Sections~\ref{f:ch:decomp} and~\ref{f:ch:regular}} answer Q2. The first proves the factor decomposition and shows why splitting by terms fails. The second asks when the pieces keep the regularity of $u$.
\item \textbf{Section~\ref{f:ch:interaction}} answers Q3: the exact interaction term, its cost, and its instability.
\item \textbf{Section~\ref{f:ch:nonlin}} extends the three questions to nonlinear equations.
\item \textbf{Section~\ref{ch:algo}} turns the theorems into an algorithm with two fully worked examples.
\item \textbf{Sections~\ref{ch:systems}--\ref{ch:dictionary}} (Part~II) go beyond constant coefficients: systems and gauge, variable coefficients, convergence and noise, the analytic encoder, and the universal dictionary.
\item \textbf{Section~\ref{ch:phase10}} records what Phase~10 changed in the statements of the earlier phases.
\item \textbf{Section~\ref{ch:extalgo}} extends the algorithm to all of the above, with a contract for every step.
\item \textbf{Section~\ref{f:ch:status}} states what is impossible, what is open, and how reliable each result is, for the whole project.
\end{itemize}

